% -----------------------------------------------------------------------------
% This file is automatically generated.
% Do not edit manually.
% Source: notebooks/support/regression-analysis/support.Rmd
% -----------------------------------------------------------------------------


\noindent A practical way to address first-order autocorrelation is to apply an AR(1) transformation. This notebook shows the transformation explicitly rather than calling a separate named routine. We first estimate the lag-1 autocorrelation in the residuals, then quasi-difference the response and the predictors, and finally refit the model to the transformed data.

\noindent We first estimate an AR(1) model on the residuals from \ttblue{mod.3}.

\par\addvspace{\topsep}
\begingroup
\fvset{listparameters={%
  \setlength{\topsep}{0pt}%
  \setlength{\partopsep}{0pt}%
  \setlength{\parsep}{0pt}%
  \setlength{\itemsep}{0pt}%
}}
\begin{Verbatim}[breaklines=true,formatcom=\color{ada_blue}]
(rho <- arima(ussteamco.mod.3$residuals,
  order = c(1, 0, 0)
)$coef[1])
\end{Verbatim}
\begin{Verbatim}[breaklines=true,formatcom=\color{ada_light_blue}]
      ar1
0.4413799
\end{Verbatim}
\endgroup
\par\addvspace{\topsep}

\noindent Then we transform all variables using the estimated \(\rho\). When the model includes interaction terms, we first construct the interaction terms in their original form and only then transform those interaction terms. We should not multiply separately transformed main effects.
